\documentclass[runningheads]{llncs}

\usepackage{listings}
\usepackage{xcolor}
\usepackage[T1]{fontenc}
\usepackage{graphicx}

\lstset{
    language=HTML,
    basicstyle=\ttfamily\small,
    numbers=left,
    numberstyle=\tiny,
    frame=single,
    breaklines=true
}

\begin{document}

\title{Informe del Módulo 1}

\author{Maia Ortiz \and Celina Lopez \and Rocio Gomez \and Bautista Pons}
\authorrunning{Ortiz et al.}
\institute{Universidad Nacional de Cuyo}

\maketitle

\begin{abstract}
El uso de Markdown permite crear páginas en GitHub de forma sencilla, estructurando contenido con sintaxis clara. Estas herramientas, junto con HTML y LaTeX, facilitan la documentación técnica y el trabajo colaborativo \cite{lopez2003latex}. Además, la construcción de páginas web y el uso de plataformas como Colab amplían las posibilidades de desarrollo \cite{sepulcre2024uso}.
\end{abstract}

\section{Introducción}

El presente documento constituye la entrega del Módulo 1, cuyo objetivo es adquirir conocimientos en lenguajes de marcado y herramientas digitales.

Se analizan HTML, Markdown, LaTeX y Google Colab como herramientas clave para la formación en ingeniería.

\section{Uso Básico de HTML}

\begin{lstlisting}
<html>
  <head>
    <title>Mi primera página web</title>
  </head>
  <body>
    <h1>Mi primera página web</h1>
    <p>Esto es un párrafo</p>
    <ul>
      <li>primero</li>
      <li>segundo</li>
    </ul>
  </body>
</html>
\end{lstlisting}

HTML es el lenguaje base para la creación de páginas web, permitiendo estructurar contenido mediante etiquetas.

\section{Lenguajes: Markdown y LaTeX}

Markdown permite generar documentos simples y claros, siendo ampliamente utilizado en GitHub.

LaTeX, en cambio, permite generar documentos académicos con gran precisión, siendo ideal para informes técnicos.

\subsection{Markdown en GitHub}

Markdown convierte texto plano en HTML estructurado, facilitando la documentación.

\subsection{Escritura en LaTeX}

LaTeX separa contenido y formato, permitiendo generar documentos profesionales con manejo automático de referencias.

\section{Cheatsheets}

Las cheatsheets permiten consultar rápidamente comandos y sintaxis, facilitando el trabajo técnico.

\section{Google Colab}

Google Colab permite ejecutar código en la nube sin necesidad de instalar software, siendo clave para análisis de datos.

\section{Conclusión}

El módulo permitió integrar herramientas fundamentales para la documentación y el trabajo técnico, sentando bases importantes para la formación profesional.

\bibliographystyle{splncs04}
\bibliography{Bibliografía}

\end{document}
